% Part: second-order-logic
% Chapter: metatheory
% Section: second-order-arithmetic

\documentclass[../../../include/open-logic-section]{subfiles}

\begin{document}

\olfileid{sol}{met}{spa}

\olsection{حساب الرتبة الثانية}

تذكّر أن نظرية حساب بيانو~$\Th{PA}$ تشمل البديهيات الثماني لـ$\Th{Q}$:
\begin{align*}
& \lforall[x][\eq/[x'][\Obj 0]]\\
& \lforall[x][\lforall[y][(\eq[x'][y'] \lif \eq[x][y])]]\\
& \lforall[x][(\eq[x][\Obj 0] \lor \lexists[y][\eq[x][y']])]\\ 
& \lforall[x][\eq[(x + \Obj 0)][x]]\\
& \lforall[x][\lforall[y][\eq[(x + y')][(x + y)']]]\\
& \lforall[x][\eq[(x \times \Obj 0)][\Obj 0]]\\
& \lforall[x][\lforall[y][\eq[(x \times y')][((x \times y) + x)]]]\\
& \lforall[x][\lforall[y][(x < y \liff \lexists[z][\eq[(z' + x)][y]])]]\\
\intertext{مضافًا إليها جميع الجمل من الصورة}
& (!A(\Obj 0) \land \lforall[x][(!A(x) \lif !A(x'))]) \lif \lforall[x][!A(x)].
\end{align*}
والأخيرة «مخطط»، أي نمط يولد عددًا لانهائيًا من~!!{sentence}s لغة
الحساب، واحدة لكل~!!{formula}~$!A(x)$. ونسمي هذا المخطط
\emph{مخطط بديهية الاستقراء} (من الرتبة الأولى).
أما في حساب بيانو \emph{من الرتبة الثانية}~$\Th{PA^2}$، فيمكن تقرير
الاستقراء بجملة واحدة. وتتألف~$\Th{PA^2}$ من البديهيات الثماني الأولى
أعلاه، مضافًا إليها \emph{بديهية الاستقراء} (من الرتبة الثانية):
\[
\lforall[X][((X(\Obj 0) \land \lforall[x][(X(x) \lif X(x'))]) \lif \lforall[x][X(x)])].
\]
وهي تقول إنه إذا كانت مجموعة جزئية~$X$ من الـ!!{domain} تحتوي
على~$\Assign{\Obj{0}}{M}$، وتحتوي مع كل~$x \in \Domain{M}$ أيضًا
على~$\Assign{\prime}{M}(x)$ (أي إنها «مغلقة تحت الخلف»)، فإنها تحتوي
كل ما في الـ!!{domain} (أي إن~$X = \Domain{M})$.

وتضمن بديهية الاستقراء ألا تحتوي أي~!!{structure} تحققها إلا
الـ!!{element}s من~$\Domain{M}$ التي تقتضي البديهيات وجودها؛ أي قيم
$\num{n}$ لكل~$n \in \Nat$. فلا يحتوي أي نموذج لـ$\Th{PA^2}$ أعدادًا
غير قياسية.

\begin{thm}
\ollabel{thm:sol-pa-standard}
إذا كانت~$\Sat{M}{\Th{PA^2}}$، فإن~$\Domain{M} =
\Setabs{\Value{\num{n}}{M}}{n \in \Nat}$.
\end{thm}

\begin{proof}
لتكن~$N = \Setabs{\Value{\num{n}}{M}}{n \in \Nat}$، وافترض أن
$\Sat{M}{\Th{PA^2}}$. ومن الواضح، لكل~$n \in \Nat$، أن
$\Value{\num{n}}{M} \in \Domain{M}$؛ ومن ثم~$N \subseteq \Domain{M}$.

ننتقل الآن إلى الاحتواء في الاتجاه الآخر. تأمل تعيين متغيرات~$s$
بحيث~$s(X) = N$. وبحسب الفرض،
\begin{align*}
\Sat{M}{\lforall[X][((X(\Obj 0) \land \lforall[x][(X(x)
      \lif X(x'))]) \lif \lforall[x][X(x)])]}, & \text{ومن ثم}\\
\Sat{M}{(X(\Obj 0) \land \lforall[x][(X(x) \lif X(x'))]) \lif
  \lforall[x][X(x)]}[s]. & 
\end{align*}
تأمل مقدم هذه الشرطية. لدينا~$\Value{\Obj{0}}{M} \in N$، ومن ثم
$\Sat{M}{X(\Obj{0})}[s]$. والحد الثاني،
$\lforall[x][(X(x) \lif X(x'))]$، محقق أيضًا. فلنفترض أن~$x \in
N$. بحسب تعريف~$N$، يكون~$x = \Value{\num{n}}{M}$ لبعض~$n$. ومن ذلك
نحصل على~$\Assign{\prime}{M}(x) = \Value{\num{n+1}}{M} \in N$. إذن
$\Assign{\prime}{M}(x) \in N$.

وهكذا لدينا~$\Sat{M}{X(\Obj 0) \land \lforall[x][(X(x) \lif X(x'))]}[s]$.
وبالتالي~$\Sat{M}{\lforall[x][X(x)]}[s]$. لكن هذا يعني أن
$x \in s(X) = N$ لكل~$x \in \Domain{M}$. ومن ثم~$\Domain{M}
\subseteq N$.
\end{proof}

\begin{cor}
\ollabel{cor:sol-pa-categorical}
أي نموذجين لـ$\Th{PA^2}$ متماثلان بنيويًا.
\end{cor}

\begin{proof}
بحسب~\olref{thm:sol-pa-standard}، تستنفد
$\Value{\num{n}}{M}$ مجال أي نموذج لـ$\Th{PA^2}$. وكل نموذج من هذا
القبيل نموذج أيضًا لـ$\Th{Q}$. وبحسب
\olref[mod][mar][stm]{prop:thq-standard}، يكون كل نموذج كهذا قياسيًا؛
أي متماثلًا بنيويًا مع~$\Struct{N}$.
\end{proof}

عرّفنا~$\Th{PA^2}$ أعلاه بأنها النظرية التي تحتوي البديهيات الحسابية
الثماني الأولى وبديهية الاستقراء من الرتبة الثانية. والواقع أنه، بفضل
القوة التعبيرية لمنطق الرتبة الثانية، لا يحتاج حساب بيانو من الرتبة
الثانية إلا إلى البديهتين الحسابيتين \emph{الأوليين} مع الاستقراء.

\begin{prop}\ollabel{prop:sol-pa-definable}
لتكن~$\Th{PA^{2\dagger}}$ نظرية الرتبة الثانية التي تحتوي البديهتين
الحسابيتين الأوليين (بديهيتي الخلف) وبديهية الاستقراء من الرتبة الثانية.
عندئذ تكون~$\le$ و$+$ و$\times$ قابلة للتعريف في
$\Th{PA^{2\dagger}}$.
\end{prop}

\begin{proof}
لإثبات أن~$\le$ قابلة للتعريف، علينا أن نجد صيغة~$!A_\le(x, y)$
بحيث تكون~$\Sat{N}{!A_\le(\num n, \num m)}$ إذا وفقط إذا كان
$n \le m$. تأمل الصيغة
\[
!B(x, Y) \ident Y(x) \land \lforall[y][(Y(y) \lif Y(y'))]
\]
من الواضح أن~$!B(\num n, Y)$ تحققها مجموعة~$Y \subseteq \Nat$ إذا
وفقط إذا كانت~$\Setabs{m}{n \le m} \subseteq Y$؛ ومن ثم يمكننا أن
نأخذ~$!A_\le(x, y) \ident
\lforall[Y][(!B(x, Y) \lif Y(y))]$.

ولرؤية أن الجمع قابل للتعريف، لاحظ أن~$k+l=m$ إذا وفقط إذا وجدت
دالة~$u$ بحيث~$u(0)=k$ و$u(n')=u(n)'$ لكل~$n$، و$m = u(l)$.
ويمكننا استعمال هذا التكافؤ لتعريف الجمع في~$\Th{PA^{2\dagger}}$
بالصيغة الآتية:
\[
!A_+(x,y,z) \ident \lexists[u][(u(\Obj 0)=x \land \lforall[w][u(w')=u
 (w)'] \land u(y)=z)]
\] 
وينبغي أن يكون واضحًا أن~$\Sat{N}{!A_+(\num k, \num l, \num m)}$
إذا وفقط إذا كان~$k+l=m$.
\end{proof}

\begin{prob}
أكمل برهان~\olref[sol][met][spa]{prop:sol-pa-definable}.
\end{prob}

\end{document}
